% exp-fips203-mlkem-pke-keygen-k2 — FIPS 203 Alg.13 ML-KEM-512 PKE KeyGen
% 编译：bash scripts/xelatex-clean.sh examples/incubating/ml-kem/ml-kem-512/exp-fips203-mlkem-pke-keygen-k2/exp-fips203-mlkem-pke-keygen-k2-实现方案-customspec.tex
\documentclass[UTF8,a4paper,11pt]{ctexart}
\usepackage{amsmath,amssymb}
\usepackage{booktabs}
\usepackage{geometry}
\usepackage{hyperref}
\usepackage{longtable}
\usepackage{array}
\usepackage{seqsplit}
\geometry{margin=2.2cm}
\newcolumntype{L}[1]{>{\raggedright\arraybackslash}p{#1}}
\newcommand{\apiname}[1]{\texttt{\seqsplit{#1}}}
\newcommand{\dirname}[1]{\texttt{\seqsplit{#1}}}
\hypersetup{colorlinks=true,linkcolor=blue,urlcolor=blue}

\title{exp-fips203-mlkem-pke-keygen-k2\\FIPS 203 Alg.13 ML-KEM-512 PKE KeyGen 实现方案}
\author{ascendc 工程（ML-KEM-512 W4a / E13）}
\date{2026-07-27}

\begin{document}
\maketitle

\begin{abstract}
本文档描述 \dirname{examples/incubating/ml-kem/ml-kem-512/exp-fips203-mlkem-pke-keygen-k2/}：
在 \textbf{Atlas A2 / Ascend910B4} 上以 AscendC 实现 FIPS 203 \textbf{Algorithm 13}，即 ML-KEM-512（$k=2$）\textbf{PKE KeyGen}。

\textbf{参数锁定}：几何、I/O 与 launch 全部锁定到
\dirname{docs/specs/fips203-mlkem512-parameter-card.md}（S-1 / T-B2）与活跃探针
\dirname{ascendc-tests/ml-kem/ml-kem-512/pass-fix-f203-alg13-device-keygen-k2/STATUS.md}。
\textbf{生产 I/O}：\texttt{input/} 仅 \texttt{seed\_d.bin} 与 NTT LUT；
\texttt{output/} 仅 \texttt{ek\_pke.bin}（800B）与 \texttt{dk\_pke.bin}（768B）。
\textbf{计划 AI Core 数}：2 launch；Launch 1 为 AIV\_ONLY \texttt{blockDim=2}，Launch 2 为 MIX \texttt{blockDim=1}（1 AIC + 2 AIV，单 cube）。
\textbf{禁止}：写入 \dirname{examples/stable/ml-kem/ml-kem-512/}；从任何 \dirname{frozen/} 目录复制、改写或移植源码；用 k3/k4 的 ek/dk 长度或零垫到 3/4/6/8 替代 k2 几何。
\end{abstract}

\tableofcontents
\newpage

\section{工程边界}

\subsection{允许的代码来源与依赖}

\begin{longtable}{@{}L{4.6cm}L{8.6cm}@{}}
\toprule
来源 & 用途 \\
\midrule
\dirname{ascendc-tests/ml-kem/ml-kem-512/pass-fix-f203-alg13-device-keygen-k2/} & 活跃 D13 k2 行为基线；可复制后自包含适配到本 exp，最终不得保留对探针目录的编译期 include/import \\
\dirname{examples/incubating/ml-kem/ml-kem-768/exp-fips203-mlkem-pke-keygen-k3/} & 活跃 k3 incubating 结构参考；仅 retarget 到 k2 锁定几何与字节长 \\
\dirname{library/shared/} & SHAKE/Keccak、CBD-$\eta=3$、公共设备侧头文件等共享库 \\
CANN / tikicpulib / CAModel & 编译、CPU 孪生与 SIM 运行环境 \\
\bottomrule
\end{longtable}

\subsection{禁止项}

\begin{itemize}
  \item 禁止从 \dirname{ascendc-tests/frozen/} 或 \dirname{examples/frozen/} 拿源码、customspec 或路线。
  \item 禁止将 \texttt{a\_hat}、\texttt{s\_hat}、\texttt{t\_hat} 等中间态作为默认 I/O 落盘；debug dump 只能显式开关启用。
  \item 禁止把 k3/k4 的 \texttt{ek/dk} 长度、polyvec6/8 或零垫当作默认。
  \item 禁止新建或写入 stable-512；本目录只作为 incubating 预研实现。
\end{itemize}

\section{数学语义（Alg.13，$k=2$，$n=256$，$q=3329$）}

\subsection{输入输出}

\begin{longtable}{@{}L{3.3cm}L{10.2cm}@{}}
\toprule
张量 / 文件 & 契约 \\
\midrule
\texttt{input/seed\_d.bin} & 4B little-endian \texttt{uint32}，由 Host 派生 FIPS 203 的 $d$；域分离串为 \texttt{exp-mlkem-f203-2s1e-k2:SEED\_D=} \\
\texttt{input/lut\_even\_stacked.bin}、\texttt{lut\_odd\_stacked.bin} & NTT Stage2 LUT；由本目录 golden 脚本生成；不是用户可见交付输入 \\
\texttt{output/ek\_pke.bin} & 800B，$\text{ByteEncode}_{12}(\hat t)$ 的 768B 后拼接 $\rho$ 的 32B \\
\texttt{output/dk\_pke.bin} & 768B，$\text{ByteEncode}_{12}(\hat s)$，其中 $\hat s$ 为前 2 个秘密多项式 \\
\bottomrule
\end{longtable}

\subsection{阶段公式}

\begin{enumerate}
  \item $G(d\Vert \texttt{byte}(2))\to(\rho,\sigma)$；$k$ 字节锁定为 2。
  \item 行 3--7：对 $(i,j)\in\{0,1\}^2$ 执行 \texttt{SampleNTT}，生成 $\hat A\in R_q^{2\times2}$，GM 形状锁定为 \texttt{[4,256] int32}。
  \item 行 8--15：用 $\eta_1=\eta_2=3$ 采样 $s\Vert e$ 共 4 个 poly；GM 形状锁定为 \texttt{[4,256]}（T-B2 polyvec4）。
  \item 行 16--20：polyvec4 NTT 后计算 $\hat t=\hat A\circ\hat s+\hat e$；InnerProduct 输出 $P\_OUT=2$，AIV 分片为 1+1。
  \item 行 21：在 compute launch 内融合生成 $ek\_pke=\text{ByteEncode}_{12}(\hat t)\Vert\rho$，同时写 $dk\_pke=\text{ByteEncode}_{12}(\hat s)$。
\end{enumerate}

\section{AscendC 实现规划}

\subsection{Launch 与 AI Core 规模}

\begin{longtable}{@{}L{2.5cm}L{3.0cm}L{7.6cm}@{}}
\toprule
Launch & 核类型 / blockDim & 数据划分与理由 \\
\midrule
1 prep & AIV\_ONLY，\texttt{blockDim=2} & $\hat A$ 共 4 poly 按 2+2 分给两个 AIV；$s\Vert e$ 共 4 poly 走 CBD-$\eta=3$ 写真实 polyvec4 GM。选择 2 AIV 是 D13 k2 探针已验锁定值。 \\
2 compute & MIX，\texttt{blockDim=1} & 1 AIC + 2 AIV（S-1 单 cube）；polyvec4 NTT；InnerProduct 输出 2 行，AIV0/AIV1 各负责 $\hat t[0]$ / $\hat t[1]$。禁止跨 cube 切 poly。 \\
\bottomrule
\end{longtable}

\subsection{GM 几何}

\begin{longtable}{@{}L{3.7cm}L{9.5cm}@{}}
\toprule
对象 & 锁定几何 \\
\midrule
$\hat A$ & \texttt{[4,256] int32}；4 个真实 poly，禁止 pad 到 9/16 \\
$s\Vert e$ & \texttt{[4,256]}；polyvec4，NTT 前后都按真实 4 poly 管理 \\
NTT S0 & \texttt{[8,256] int8}，即 \texttt{[HI4, LO4]}；每个 AIV 持有完整 poly 的 hi+lo，禁止 limbsplit \\
Stage2 mat\_c & \texttt{[32,128] int32}；Cube 的 M 对齐垫到 16 只是硬件对齐，不表示假 poly \\
ByteEncode12 & $2\times384=768$B；分别用于 $\hat t$ 与 $\hat s$；ek 再拼 $\rho$ 得 800B \\
\bottomrule
\end{longtable}

\subsection{同步与流水}

Prep 内部使用 UB/GM 分段搬运与必要的 \apiname{PipeBarrier}；两次 launch 之间由 Host 顺序保证 GM 可见性。
Compute 内部遵守 B5/B6/D13 已验模式：AIV 写入平面 Stage1/Stage3 GM 后以 CrossCore flag 与 AIC 同步；AIC 写完 \texttt{mat\_c} 后 AIV 等待再做 Stage3 merge、InnerProduct 与 ByteEncode12。
CPU 仅作为逻辑孪生；SIM 是同步与搬运验收门槛。

\section{AscendC API 表}

本实现不引入 E13 独有的新 AscendC API；下表 API 均已有索引记录，使用约束按
\dirname{library/documents/CANN-AscendC算子开发接口参考-查阅索引.md} 执行。

\begin{longtable}{@{}L{3.2cm}L{2.7cm}L{2.5cm}L{2.2cm}L{3.0cm}@{}}
\toprule
API 名 & 分类 / 文档 & Atlas A2 & 本用例 dtype & 备注 \\
\midrule
\apiname{GetBlockIdx} / \apiname{GetSubBlockIdx} & 系统变量 / 资源 & √ & 标量索引 & 区分 AIV 分片与 MIX 子核 \\
\apiname{GlobalTensor} / \apiname{LocalTensor} & 基础结构 & √ & \texttt{uint8/int8/int32} & GM/UB 张量契约须与本 spec 几何一致 \\
\apiname{DataCopy} & Memory 数据搬运 & √ & \texttt{uint8/int8/int32} & 连续块搬运；32B 对齐约束 \\
\apiname{Duplicate} & Memory 矢量计算 & √ & \texttt{uint8/int32} & UB 清零、pad 与 tail buffer 初始化 \\
\apiname{PipeBarrier} & Pipe 同步 & √ & — & MTE/Vector/Cube 阶段交界显式同步 \\
\apiname{CrossCoreSetFlag} / \apiname{CrossCoreWaitFlag} & MIX 同步 & √ & — & AIV/AIC 通过 GM 交接；CPU 过不能删 \\
\apiname{Mmad} / \apiname{Fixpipe} & 矩阵计算 & √ & \texttt{int8*int8->int32} & Stage2 MMAD；硬件 M 对齐垫不表示假 poly \\
\apiname{ShiftRight} & 矢量算术 & √ & \texttt{int32} & limb 拆分、Barrett 与 ByteEncode 辅助右移 \\
\apiname{Muls} / \apiname{Add} / \apiname{Sub} & 矢量算术 & √ & \texttt{int32/int16} & limb、模加减、压缩/编码辅助 \\
\apiname{Cast} & 类型转换 & √（受限） & \texttt{int32/int16/half/int8} & 遵守 A2 Cast 链；不得假设 \texttt{int32->int8} 单跳 \\
\apiname{Gather} & 矢量搬运/选择 & √ & \texttt{int32} & 仅允许非 NTT S1--S3 段；NTT 三段内禁用 \\
\apiname{GetValue} / \apiname{SetValue} & LocalTensor 标量访问 & √ & \texttt{uint8/int32} & 仅用于 pack 尾部等标量缺口 \\
\bottomrule
\end{longtable}

\subsection{评估过但未采用}

\begin{longtable}{@{}L{3.8cm}L{9.2cm}@{}}
\toprule
API / 路线 & 不采用原因 \\
\midrule
\apiname{Compares}/\apiname{GatherMask} compact 链 & Alg.13 E13 不需要新增 reject compact；沿用 D13 已验 \texttt{Mins}/标量 rej 路径 \\
高阶 \apiname{Matmul<>} & 继承 D13 手写 MMAD + 平面 mat\_c 契约 \\
\apiname{Gather} 用于 NTT S1--S3 & 参数卡与 ML-KEM NTT 定稿禁止 \\
零垫到 polyvec6/8 & S-1/T-B2 锁定禁零垫 \\
\bottomrule
\end{longtable}

\section{数学步骤覆盖率}

\begin{longtable}{@{}L{4.2cm}L{4.0cm}L{2.0cm}L{3.0cm}@{}}
\toprule
数学步骤 & 实现方式 / API & 覆盖 & 缺口说明 \\
\midrule
$G(d\Vert2)$ 与 XOF & shared Keccak/SHAKE + UB 流水 & 100\% 设备 & 标量 Keccak 核属本仓 shared \\
\texttt{SampleNTT} 生成 $\hat A[4]$ & AIV\_ONLY + rej 路径 + \apiname{DataCopy} & 100\% 设备 & 2+2 分片；无零垫 \\
CBD $\eta=3$ 生成 $s\Vert e[4]$ & AIV\_ONLY + CBD-$\eta=3$ & 100\% 设备 & 输出按 polyvec4 真实行数 \\
polyvec4 NTT S1--S3 & MIX AIV/AIC + \apiname{Mmad} & 100\% 设备 & NTT 内无 \apiname{Gather} \\
InnerProduct $P=2$ & Alg.11 向量路径 + AIV 1+1 & 100\% 设备 & 整除分片 \\
ByteEncode12 + $ek\Vert\rho$ & ByteEncode12 向量 + 尾部标量 pack & 部分向量 & pack 尾部标量缺口沿用 D13 已验实现 \\
\bottomrule
\end{longtable}

\section{验证计划}

\subsection{默认验收}

\begin{verbatim}
cd examples/incubating/ml-kem/ml-kem-512/exp-fips203-mlkem-pke-keygen-k2
bash run.sh -r cpu -v Ascend910B4
SIM_DIRECT=1 bash run.sh -r sim -v Ascend910B4
\end{verbatim}

期望：\texttt{ek\_pke.bin} 为 800B，\texttt{dk\_pke.bin} 为 768B；若开启 \texttt{KEYGEN\_VERIFY=1}，两者均与 golden 对拍 PASS。
SIM 通过后检查用例根目录无 stray \texttt{core*.dump}、\texttt{profile\_*log*.toml}。

\subsection{可选对照}

后续 glue 补 liboqs-512 KAT；E13 本轮门禁为 CPU + \texttt{SIM\_DIRECT=1} sim。
liboqs 仅作为黑盒 I/O oracle；禁止把 liboqs 源码实现同构移植进 AscendC。

\section{产出物}

\begin{itemize}
  \item \dirname{exp-fips203-mlkem-pke-keygen-k2-实现方案-customspec.tex}
  \item \dirname{exp-fips203-mlkem-pke-keygen-k2-实现方案-customspec.pdf}
  \item 后续实现文件须与本 customspec 同目录自包含，且自研 Python/C/AscendC 写详细中文注释。
\end{itemize}

\end{document}
